\documentclass[11pt]{article}
\usepackage[T1]{fontenc}
\usepackage{lmodern}
\usepackage[margin=1in]{geometry}
\usepackage{amsmath,amssymb,amsthm,mathtools,booktabs}
\usepackage{microtype}
\usepackage[hidelinks]{hyperref}
\hypersetup{pdftitle={},pdfauthor={},pdfsubject={},pdfkeywords={}}
\newcommand{\tr}{\operatorname{tr}}
\newcommand{\Hess}{\operatorname{Hess}}
\newcommand{\dd}{\mathrm d}
\newcommand{\ii}{\mathrm i}
\newcommand{\Gden}{\mathcal G}
\newtheorem{proposition}{Proposition}
\newtheorem{remark}{Remark}
\title{Native Mixed Quartic Interactions on an Einstein--Scalar Relative Carrier\\
\large and an Exact Matter-Supported Walker Geometry}
\author{}
\date{29 October 2026}
\begin{document}
\maketitle
\begin{abstract}
We construct the native two-graviton/two-scalar quartic contact of a flat Einstein--massless-scalar theory after eliminating its light-front auxiliary metric variables.
The construction keeps arbitrary off-shell external momenta and gives a finite tensor formula on a curvature-normalized two-helicity relative carrier.
The scalar action supplies an explicit intrinsic mixed contact, while the Einstein action supplies a quadratic gravitational source and its auxiliary feedback.
The latter source contains an external-wave correction, proportional to the sum of the two graviton wave symbols, that is absent from the conserved on-shell current.
We derive the correction, state every source and Hessian, and independently check its normalization against the Einstein Gamma--Gamma action.
The contact has only the declared longitudinal chart poles and no new inverse wave symbol.
As a distinct finite-sector construction, we solve the full Einstein--scalar equations in a Walker family with exact scalar backreaction and a rank-two scalar Hessian.
That invariant distinguishes the solution from an aligned-scalar pp wave.
Its real section has split signature; it is not a Lorentzian Real lift of the localized carrier.
Quartic chart descent, the full gauge--antifield structure and the quantum measure remain unconstructed.
\end{abstract}

\section{Problem and conventions}

A contact generated by a quadratic field redefinition is not, by itself, the native quartic interaction of the original action.
Likewise, a conserved two-graviton current evaluated on free shells does not supply the off-shell source needed to eliminate auxiliary variables inside a higher interaction.
We address these issues directly in one localized Einstein--scalar carrier and construct a complementary finite matter-supported geometry.

The starting action is
\begin{equation}
 S[g,\phi]=\frac{1}{2\kappa}\int\dd^4x\,\sqrt{|g|}\,R
 -\frac12\int\dd^4x\,\sqrt{|g|}\,g^{\mu\nu}\partial_\mu\phi\partial_\nu\phi,
 \qquad \kappa>0.
 \label{eq:action}
\end{equation}
There is one minimally coupled real massless scalar and no scalar potential.
Our curvature convention is
\begin{equation}
 R_{\mu\nu}=\partial_\rho\Gamma^\rho_{\mu\nu}
 -\partial_\nu\Gamma^\rho_{\mu\rho}
 +\Gamma^\rho_{\rho\lambda}\Gamma^\lambda_{\mu\nu}
 -\Gamma^\rho_{\nu\lambda}\Gamma^\lambda_{\mu\rho}.
 \label{eq:ricci-convention}
\end{equation}
For bulk action coefficients we use the integrated Einstein Gamma--Gamma representative.
Its difference from the Einstein--Hilbert volume density is a variational surface term; no null-boundary action, asymptotic prescription or measure is assumed.

The flat calculation uses
\begin{equation}
 \eta=-2\dd u\,\dd v+\dd x^2+\dd y^2,\qquad
 p=(d,a,p_x,p_y)^T,\qquad
 W(p)=-p^T\eta p=2ad-p_x^2-p_y^2.
 \label{eq:flat}
\end{equation}
Momenta are covectors, and a Fourier derivative is \(\partial_\mu\mapsto\ii p_\mu\).
The matrices of \(\eta\) and its inverse agree in this basis.
All four-point kernels have incoming momenta summing to zero, with metric legs 1,2 and scalar legs 3,4.
No external wave symbol is set to zero in the main quartic construction.

The interpretation of auxiliary elimination as tree-level homotopy transfer is established in Ref.~\cite{Arvanitakis2020}.
Here the required quantities are calculated explicitly rather than inferred from that general framework.
The finite construction uses a heavenly-type Walker equation; the established heavenly framework is discussed, for example, in Ref.~\cite{Krasnov2024}.
Neither reference substitutes for the explicit off-shell normalization or finite matter-supported solution below.

\section{The physical inclusion and its relative coordinates}
\label{sec:inclusion}

Write the metric in a fixed light-front gauge as
\begin{equation}
 g_{v\mu}\dd x^\mu=-\dd u,\qquad
 g_{uu}=A,\qquad g_{ui}=B_i,\qquad g_{ij}=T_{ij},
 \quad i,j\in\{x,y\}.
 \label{eq:lfmetric}
\end{equation}
The physical transverse field is trace free:
\begin{equation}
 C=\begin{pmatrix}f&g\\g&-f\end{pmatrix},\qquad
 T=I+C,\qquad \rho^2=f^2+g^2,\qquad C^2=\rho^2 I.
 \label{eq:C}
\end{equation}
The symbol \(g\) inside \(C\) denotes one transverse coefficient; metric components retain indices.
The linear physical inclusion fixes its longitudinal tails by
\begin{equation}
 B_i=\partial_v^{-1}\partial_jC_{ij},\qquad
 A=\partial_v^{-2}\partial_i\partial_jC_{ij}.
 \label{eq:tails}
\end{equation}
These are rational Fourier multipliers on the nonzero longitudinal chart; no retarded inverse is identified with a cyclic inverse.
For a single covector,
\begin{equation}
 B=\frac{Cp_\perp}{a},\quad
 A=\frac{p_\perp^TCp_\perp}{a^2},\quad
 h(p,C)=\begin{pmatrix}
 A&0&B_x&B_y\\0&0&0&0\\B_x&0&f&g\\B_y&0&g&-f
 \end{pmatrix}.
 \label{eq:physical-h}
\end{equation}
It follows directly, including off shell, that
\begin{equation}
 \tr(\eta h)=0,\qquad p^T\eta h=0.
 \label{eq:TT}
\end{equation}
Canonical metric fields are inserted by \((f,g)=\sqrt{2\kappa}(F,G)\).
Until that substitution is made, these are literal metric coefficients.

The curvature-normalized relative carrier can be defined without assuming a nonlinear twistor correspondence.
Use complex transverse coordinates with \(\eta=-2\dd u\,\dd v+2\dd z\,\dd w\), and write \(p=(d,a,b,c)^T\).
Set
\begin{equation}
 e_z=(c/a,0,1,0)^T,\qquad e_w=(b/a,0,0,1)^T,
 \qquad \Delta=ad-bc,\qquad W=2\Delta.
 \label{eq:ezew}
\end{equation}
The literal relative fields \(r_B,t_B\) have metric inclusion
\begin{equation}
 \boxed{h_B(p;r_B,t_B)=\kappa a^2r_B\,e_we_w^T
                         +\frac{2}{a^2}t_B\,e_ze_z^T.}
 \label{eq:anchor}
\end{equation}
Its linearized Weyl components on the radiative shell are
\begin{equation}
 C^{(1)}_{vzvz}=t_B,\qquad
 C^{(1)}_{vwvw}=\frac{\kappa a^4}{2}r_B.
 \label{eq:weyl-anchor}
\end{equation}
These fix the reciprocal pivot factors in the curvature anchor.
The free relative Hessian, including the scalar, is
\begin{equation}
 K(p)=\Delta J,\qquad
 J=\begin{pmatrix}0&1&0\\1&0&0\\0&0&2\end{pmatrix},\qquad
 x=(r_B,t_B,\phi)^T.
 \label{eq:relative-K}
\end{equation}
The scalar factor two is its massless wave normalization.
Equation~\eqref{eq:anchor} specifies the literal carrier input to the quartic calculation, localized to one chart.

\section{The native mixed quartic from the scalar action}
\label{sec:native}

The determinant and inverse of Eq.~\eqref{eq:lfmetric} give
\begin{equation}
 \sqrt{|g|}=\nu=\sqrt{\det T}=\sqrt{1-\rho^2},\qquad
 g^{uv}=-1,\quad g^{ij}=(T^{-1})_{ij},\quad g^{vi}=(T^{-1}B)_i,
 \label{eq:inverse1}
\end{equation}
and \(g^{vv}=-A+B^TT^{-1}B\).
The exact scalar density is
\begin{align}
 \mathcal L_\phi=\nu\bigg[
 &\phi_u\phi_v+\frac{A}{2}\phi_v^2
 -\phi_vB^TT^{-1}\nabla_\perp\phi
 -\frac12\nabla_\perp\phi^TT^{-1}\nabla_\perp\phi\notag\\
 &-\frac12B^TT^{-1}B\,\phi_v^2\bigg].
 \label{eq:scalar-exact}
\end{align}
This permits arbitrary transverse dependence.
On the linear physical inclusion, \(A,B,C\) each have metric degree one.
Using
\begin{equation}
 \nu=1-\frac{\rho^2}{2}+O(C^4),\qquad
 \nu T^{-1}=I-C+\frac{\rho^2}{2}I+O(C^3)
 \label{eq:nu-expansion}
\end{equation}
gives the native mixed quartic density
\begin{equation}
 \boxed{\mathcal L^{\rm native}_{C^2\phi^2}
 =-\frac{\rho^2}{2}\phi_u\phi_v
 +\phi_vB^TC\nabla_\perp\phi
 -\frac{\rho^2}{4}|\nabla_\perp\phi|^2
 -\frac{|B|^2}{2}\phi_v^2.}
 \label{eq:native-density}
\end{equation}
There is no \(AC\phi^2\) term: the exact inverse is affine in \(A\), and the first two degrees of \(\nu A\) are one and three.
The terms containing \(B\) are essential for transverse-dependent states.
Keeping only the transverse block would miss native interactions before any auxiliary elimination.

For two TT metric inclusions define
\begin{equation}
 D_{12}=\eta h_1\eta h_2\eta+\eta h_2\eta h_1\eta
 -\frac12\tr(\eta h_1\eta h_2)\eta.
 \label{eq:D12}
\end{equation}
This is the mixed coefficient in \(\sqrt{|g|}g^{-1}\), since both linear traces vanish.
The scalar derivatives contribute \(\ii^2=-1\), canceling the action sign, and the two assignments of scalar labels cancel the factor \(1/2\).
The polarized native contact is consequently
\begin{equation}
 \boxed{Q^{\rm native}_{12;34}=p_3^TD_{12}p_4.}
 \label{eq:native-contact}
\end{equation}
Equations~\eqref{eq:native-density} and \eqref{eq:native-contact} express the same intrinsic action term.
Neither is inferred solely from a nonlinear coordinate change.

\section{The actual quadratic Einstein source off shell}
\label{sec:source}

Gravitational auxiliary feedback requires the quadratic source obtained by varying the Einstein action in auxiliary metric directions.
An on-shell conserved current is insufficient.

\subsection{Finite tensor formula for the mixed Ricci coefficient}

Define the Fourier connections
\begin{align}
 \Gamma_i{}^\rho_{\mu\nu}
 &=\frac{\ii}{2}\eta^{\rho\sigma}
 (p_{i\mu}h_{i\sigma\nu}+p_{i\nu}h_{i\sigma\mu}-p_{i\sigma}h_{i\mu\nu}),
 \label{eq:gamma1}\\
 \Gamma_{12}{}^\rho_{\mu\nu}
 &=-(\eta h_1)^\rho{}_{\sigma}\Gamma_2{}^\sigma_{\mu\nu}
   -(\eta h_2)^\rho{}_{\sigma}\Gamma_1{}^\sigma_{\mu\nu}.
 \label{eq:gamma12}
\end{align}
For \(k=p_1+p_2\),
\begin{align}
 R^{12}_{\mu\nu}={}&\ii k_\rho\Gamma_{12}{}^\rho_{\mu\nu}
 -\ii k_\nu\Gamma_{12}{}^\rho_{\mu\rho}\notag\\
 &+\Gamma_1{}^\rho_{\rho\lambda}\Gamma_2{}^\lambda_{\mu\nu}
  +\Gamma_2{}^\rho_{\rho\lambda}\Gamma_1{}^\lambda_{\mu\nu}\notag\\
 &-\Gamma_1{}^\rho_{\nu\lambda}\Gamma_2{}^\lambda_{\mu\rho}
  -\Gamma_2{}^\rho_{\nu\lambda}\Gamma_1{}^\lambda_{\mu\rho}.
 \label{eq:ric12}
\end{align}
This is an explicit finite formula in the two momenta and physical metric coefficients.
It imposes neither a free equation nor conservation of the quadratic current, and uses only the longitudinal chart denominators of Eq.~\eqref{eq:physical-h}.

\subsection{The external-wave correction}

For TT inclusions,
\begin{equation}
 R^{(1)}_{i,\mu\nu}=-\frac{W_i}{2}h_{i,\mu\nu},\qquad R^{(1)}_i=0.
 \label{eq:Rlinear}
\end{equation}
The vanishing linear traces mean that the volume factor makes no additional degree-two contribution to the raised Einstein density.
The inverse-metric corrections nevertheless remain.
The mixed scalar curvature is
\begin{equation}
 R^{[12]}=\tr(\eta R^{12})
 +\frac{W_1+W_2}{2}\tr(\eta h_1\eta h_2).
 \label{eq:Rscalar12}
\end{equation}
Raising both Ricci indices also produces the two ordered products of the inclusions.
Combining these terms gives the actual Euler tensor density
\begin{equation}
 \boxed{\Gden^{12}:=[\sqrt{|g|}G^{\mu\nu}]_{12}
 =\eta R^{12}\eta-\frac{\eta}{2}\tr(\eta R^{12})
 +\frac{W_1+W_2}{2}D_{12}.}
 \label{eq:density-source}
\end{equation}
The last term is generically nonzero off shell.
It vanishes on the two graviton free shells, explaining its absence from the conserved on-shell source.
Deleting it before reduction changes the off-shell quartic.

\subsection{All auxiliary source coordinates}

Separate the physical inclusion from the remaining auxiliary variables
\begin{equation}
 y=(t,A_{\rm aux},B_{x,\rm aux},B_{y,\rm aux})^T.
 \label{eq:y}
\end{equation}
The trace coordinate \(t\), distinct from the helicity variable \(t_B\), is normalized by \(\delta T=(t/2)I\).
The metric basis \(E_t\) has entries \((E_t)_{xx}=(E_t)_{yy}=1/2\), while \((E_A)_{uu}=1\).
For each \(i=x,y\), the basis \(E_{B_i}\) has \((u,i)\) and \((i,u)\) entries equal to one; all other basis entries vanish.
Since
\begin{equation}
 \delta S_{\rm EH}=-\frac{1}{2\kappa}\int
 \sqrt{|g|}G^{\mu\nu}\delta g_{\mu\nu},
 \label{eq:variation}
\end{equation}
the two-graviton source is
\begin{equation}
 \boxed{J^{CC}_{12}=-\frac{1}{2\kappa}
 \begin{pmatrix}
 (\Gden^{xx}+\Gden^{yy})/2\\\Gden^{uu}\\2\Gden^{ux}\\2\Gden^{uy}
 \end{pmatrix}_{12}.}
 \label{eq:JCC}
\end{equation}
This specifies all factors associated with the trace and off-diagonal shifts.
For scalar covectors \(p=(d,a,p_x,p_y)\), \(q=(e,b,q_x,q_y)\), direct scalar variation gives
\begin{equation}
 J^{\phi\phi}_{pq}=\begin{pmatrix}
 -(db+ae)/2\\-ab\\a q_x+b p_x\\a q_y+b p_y
 \end{pmatrix}.
 \label{eq:Jscalar}
\end{equation}
Both scalar legs may be off shell.
Around zero scalar background, no auxiliary source of degree \(C\phi\) exists because the scalar action is quadratic in \(\phi\).

\section{Auxiliary elimination and the mixed quartic}
\label{sec:quartic}

For channel covector \(k=(d,a,k_x,k_y)\), the quadratic Einstein Hessian in the declared basis is
\begin{equation}
 H(k)=\frac{1}{4\kappa}\begin{pmatrix}
 -ad&-a^2&ak_x&ak_y\\-a^2&0&0&0\\ak_x&0&2a^2&0\\ak_y&0&0&2a^2
 \end{pmatrix}.
 \label{eq:H}
\end{equation}
Its inverse is
\begin{equation}
 H^{-1}(k)=\kappa\begin{pmatrix}
 0&-4/a^2&0&0\\
 -4/a^2&4d/a^3+2k_\perp^2/a^4&2k_x/a^3&2k_y/a^3\\
 0&2k_x/a^3&2/a^2&0\\0&2k_y/a^3&0&2/a^2
 \end{pmatrix}.
 \label{eq:Hinv}
\end{equation}
The physical inclusion is orthogonal to the auxiliary directions in the quadratic action, so the auxiliary response begins at physical degree two.
The terms relevant through degree four are
\begin{align}
 S={}&S_{\rm phys,2}+S_{\rm phys,3}+S^{\rm native}_{\rm phys,4}
 +\frac12\langle y,Hy\rangle\notag\\
 &+\langle y,J^{CC}+J^{\phi\phi}\rangle+O({\rm degree}\ 5).
 \label{eq:elimination-action}
\end{align}
Thus \(y_2=-H^{-1}(J^{CC}+J^{\phi\phi})\).
Higher auxiliary corrections cannot contribute at degree four because they multiply the already stationary quadratic-plus-source equation.
Elimination produces
\begin{equation}
 S_{4,\rm aux}=-\frac12\left\langle J^{CC}+J^{\phi\phi},
 H^{-1}(J^{CC}+J^{\phi\phi})\right\rangle.
 \label{eq:schur}
\end{equation}
The two assignments of source types remove the factor \(1/2\) from the polarized mixed term.

\begin{proposition}
On the chart \(a_1a_2(a_1+a_2)\ne0\), the full mixed quartic contact of the reduced physical Einstein--scalar action is
\begin{equation}
 \boxed{Q^{\rm red}_{CC\phi\phi}(1,2;3,4)
 =p_3^TD_{12}p_4
 -(J^{CC}_{12})^TH^{-1}(p_1+p_2)J^{\phi\phi}_{34}.}
 \label{eq:main}
\end{equation}
All four external momenta may be off shell.
The contact introduces no inverse external or channel wave symbol.
\end{proposition}

\begin{proof}
The first term is the native scalar contact from Section~\ref{sec:native}.
The Einstein action itself has no scalar-dependent native term.
The only additional fourth-degree term with two scalar and two metric inputs is the mixed term in Eq.~\eqref{eq:schur}, evaluated with the actual sources of Section~\ref{sec:source}.
There is no quadratic auxiliary \(C\phi\) source creating another pairing.
Every denominator comes from the physical inclusion or the explicit inverse Hessian, hence is a power of a declared longitudinal pivot rather than a wave symbol.
\end{proof}

Scalar and radiative-graviton exchange are separate diagrams generated by this reduced action and its cubic vertices.
They are not part of the contact \eqref{eq:main} and retain their physical channel poles.
The complete on-shell Compton amplitude is therefore a different object and is not counted here as a new result.

\subsection{A closed source and actual nonlinear trace response}

The exact identity
\begin{equation}
 R_{vv}=-\frac12\tr(T^{-1}T_{vv})
 +\frac14\tr(T^{-1}T_vT^{-1}T_v)
 \label{eq:Rvv}
\end{equation}
is independent of the other metric blocks and of transverse derivatives.
For trace-free \(C\), its quadratic part is
\begin{equation}
 R_{vv}^{(2)}=\frac12\tr(CC_{vv})+\frac14\tr(C_v^2).
 \label{eq:Rvv2}
\end{equation}
Polarization gives \(R^{12}_{vv}=-(a_1^2+a_2^2+a_1a_2)\tr(C_1C_2)/2\).
Since \(\Gden^{uu}=R_{vv}\),
\begin{align}
 J^{CC}_{A,12}&=\frac{a_1^2+a_2^2+a_1a_2}{4\kappa}\tr(C_1C_2),
 \label{eq:JAclosed}\\
 t^{CC}_{12}&=\frac{a_1^2+a_2^2+a_1a_2}{(a_1+a_2)^2}\tr(C_1C_2).
 \label{eq:tclosed}
\end{align}
The second expression is the actual auxiliary area/trace response.
The scalar-sourced trace in this normalization is
\begin{equation}
 t^{\phi\phi}_{34}=-\frac{4\kappa a_3a_4}{(a_3+a_4)^2}.
 \label{eq:scalar-trace}
\end{equation}
These responses connect the transferred contact with nonlinear metric constraints.

\subsection{Literal carrier evaluation and higher chart comparison}

To evaluate Eq.~\eqref{eq:main} on the relative fields, insert Eq.~\eqref{eq:anchor} and use the complex transverse metric and transformed auxiliary basis.
Equivalently, transform Eq.~\eqref{eq:anchor} into the real transverse basis before evaluating the real-coordinate formulas.
This supplies every two-helicity choice \((r_B,r_B)\), \((r_B,t_B)\) and \((t_B,t_B)\) with two scalar inputs.

A cubic-compatible chart map does not automatically identify the native quartic in another chart.
For a free isometry \(L\) and a field map \(f=L(x+u_2+u_3+\cdots)\), the next defect is
\begin{align}
 D_4={}&S_{4,b}(Lx)-S_{4,c}(x)
 +\dd S_{3,b}(Lx)[Lu_2]
 +\frac12\langle u_2,Ku_2\rangle.
 \label{eq:quartic-defect}
\end{align}
Solving for \(u_3\) requires decomposing this defect in the ideal generated by external wave symbols, with actual native quartics present.
An off-shell representative mixing scalar and metric fields can import pure-graviton quartics into a mixed coefficient.
Those pure-graviton coefficients are not supplied here.
The present result is the mixed native input for a higher chart comparison, not a full nonlinear Real atlas or interacting Batalin--Vilkovisky (BV) equivalence.

\section{An exact finite matter-supported Walker geometry}
\label{sec:finite}

The same full equations admit a finite construction in a different analytic sector.
Use coordinates \((u,v,z,w)\) and
\begin{equation}
 g_\theta=-2\dd u\,\dd v+2\dd z\,\dd w
 +\theta_{ww}\dd u^2+2\theta_{vw}\dd u\,\dd z+\theta_{vv}\dd z^2.
 \label{eq:walker}
\end{equation}
The real coordinate section has signature \((2,2)\); a Lorentzian Real section is not assumed.
Define
\begin{equation}
 \Box_\eta=-2\partial_u\partial_v+2\partial_z\partial_w,
 \qquad E(\theta)=\Box_\eta\theta-\det\Hess_{v,w}\theta.
 \label{eq:heavenly}
\end{equation}
The finite inverse is affine in the metric perturbation, and the only Ricci components are
\begin{equation}
 R_{uu}=-E_{ww}/2,\qquad R_{uz}=-E_{vw}/2,\qquad R_{zz}=-E_{vv}/2.
 \label{eq:walker-ricci}
\end{equation}
For an aligned scalar \(\Psi(u,z)\), its gradient is null and its exact wave equation holds.
In particular, \(\Psi=\epsilon u\) requires
\begin{equation}
 E(\theta)=-\kappa\epsilon^2w^2,
 \label{eq:sourced-E}
\end{equation}
up to terms affine in \((v,w)\) annihilated by the displayed Hessian.

\subsection{The cubic sign obstruction and a real repair}

A tempting cubic ansatz \(\theta=q w^2v\) has
\begin{equation}
 \det\Hess_{v,w}\theta=-4q^2w^2.
 \label{eq:cubic-obstruction}
\end{equation}
It has zero flat wave operator, so Eq.~\eqref{eq:sourced-E} would require \(q^2=-\kappa\epsilon^2/4\).
This is a failure of the real cubic ansatz rather than an obstruction to real backreaction.
Instead write
\begin{equation}
 \theta=f(w)+\frac12h(w)v^2,\qquad
 \det\Hess_{v,w}\theta
 =h f''+\left(\frac{hh''}{2}-(h')^2\right)v^2.
 \label{eq:detrepair}
\end{equation}
The \(v^2\) term vanishes for \(h=1/H\), \(H=\alpha w+\beta\).
Taking \(f''=\kappa\epsilon^2w^2H\) repairs the source sign with real coefficients.
Discarding affine potential terms gives
\begin{equation}
 \boxed{\theta=\kappa\epsilon^2\left(\frac{\alpha w^5}{20}
 +\frac{\beta w^4}{12}\right)+\frac{v^2}{2H},\quad
 \Psi=\epsilon u,\quad H=\alpha w+\beta\ne0.}
 \label{eq:finite-theta}
\end{equation}
Its flat wave operator vanishes since it is independent of \(u,z\), and the determinant in Eq.~\eqref{eq:detrepair} is \(\kappa\epsilon^2w^2\).
The metric is
\begin{equation}
 \boxed{\begin{aligned}
 g={}&-2\dd u\,\dd v+2\dd z\,\dd w
 +\left(\frac{\alpha^2v^2}{H^3}+\kappa\epsilon^2w^2H\right)\dd u^2\\
 &-\frac{2\alpha v}{H^2}\dd u\,\dd z+\frac1H\dd z^2.
 \end{aligned}}
 \label{eq:finite-metric}
\end{equation}
It is a regular local patch for \(H\ne0\); no global extension or classification of \(H=0\) is asserted.

\begin{proposition}
The metric \eqref{eq:finite-metric} and scalar \(\Psi=\epsilon u\) solve the full minimal Einstein--scalar equations, with exact matter-supported increment
\begin{equation}
 g(\epsilon)-g(0)=\kappa\epsilon^2w^2H\,\dd u^2.
 \label{eq:backreaction}
\end{equation}
There is no uncomputed higher power of \(\epsilon\) in this solution.
\end{proposition}

\begin{proof}
Set \(P=\theta_{ww}\), \(Q=\theta_{vw}\), \(R=\theta_{vv}\).
The metric and inverse are
\begin{equation}
 g=\begin{pmatrix}P&-1&Q&0\\-1&0&0&0\\Q&0&R&1\\0&0&1&0\end{pmatrix},\quad
 g^{-1}=\begin{pmatrix}0&-1&0&0\\-1&-P&0&Q\\0&0&0&1\\0&Q&1&-R\end{pmatrix},\quad
 \det g=1.
 \label{eq:finite-inverse}
\end{equation}
The raised scalar gradient is \(-\epsilon\partial_v\), which is null and divergence free, proving the scalar equation directly.
Equation~\eqref{eq:walker-ricci} with the sourced potential gives
\begin{equation}
 R_{\mu\nu}=\kappa\epsilon^2(\dd u)_\mu(\dd u)_\nu,\qquad R=0.
 \label{eq:finite-einstein}
\end{equation}
This is precisely the Einstein equation for the null scalar stress.
Equation~\eqref{eq:backreaction} follows from the explicit metric, without an amplitude expansion.
An independent direct computation of all four-dimensional connection and Ricci components verifies the same equations without using the reduced identity \eqref{eq:walker-ricci}.
\end{proof}

\subsection{A coordinate-invariant non-pp witness}

The scalar Hessian has the nonzero block
\begin{equation}
 (\nabla\dd\Psi)\big|_{u,z}
 =\epsilon\begin{pmatrix}
 -\alpha^2v/H^3&\alpha/(2H^2)\\\alpha/(2H^2)&0
 \end{pmatrix},\qquad
 \det=-\frac{\epsilon^2\alpha^2}{4H^4}.
 \label{eq:hessian-rank}
\end{equation}
Its rank is two for \(\alpha\epsilon\ne0\).
In a pp metric with parallel null covector \(\dd U\) and aligned scalar \(\Psi=F(U)\),
\begin{equation}
 \nabla\dd\Psi=F''(U)\dd U\otimes\dd U,
 \label{eq:pp-hessian}
\end{equation}
whose rank is at most one.
Tensor rank is invariant under diffeomorphisms, so the coupled solution is not a coordinate rewriting of that aligned-scalar pp family.
The distinction concerns the pair \((g,\Psi)\), without requiring a classification of the metric alone.
The null scalar gradient remains geodesic, since \((\nabla\dd\Psi)\nabla\Psi=0\).

\begin{remark}
The finite geometry and the localized quartic have different domains.
The displayed real section has split signature, and no Lorentzian conjugate-helicity Real condition is supplied.
Its scalar and polynomial longitudinal dependence are background or zero-longitudinal data outside the \(a\ne0\) carrier of Eq.~\eqref{eq:main}.
Thus this solution does not itself extend that carrier through zero momentum.
Both the Einstein--Hilbert density and the aligned scalar kinetic density vanish when restricted to the Walker ansatz.
The equations above are the full equations varied before restriction; a potential multiplier action cannot be identified with the restricted native action or used to infer a BV equivalence.
\end{remark}

\section{Exact verification and reproducible coefficients}

The accompanying Python calculations use SymPy 1.14.0 and exact rational or symbolic arithmetic.
For four off-shell realizations, the mixed-quartic script compares all four auxiliary source entries with independently expanded coefficients of the Einstein Gamma--Gamma action.
The second computation retains the linear volume contribution of the trace auxiliary leg and does not use Eq.~\eqref{eq:density-source}.
These 16 comparisons test signs and factors invisible on free shells.

The first realization has
\begin{align}
 p_1&=(2,3,1,4),&p_2&=(-4,3,5,0),\notag\\
 p_3&=(1,-3,2,5),&p_4&=(1,-3,-8,-9),
 \label{eq:example-p}
\end{align}
with \((f_1,g_1)=(2/3,-1/4)\) and \((f_2,g_2)=(1/5,3/2)\).
Its wave symbols are
\begin{equation}
 (W_1,W_2,W_3,W_4)=(-5,-49,-35,-151).
 \label{eq:example-w}
\end{equation}
In literal metric normalization, at \(\kappa=1\),
\begin{align}
 J^{CC}_{12}&=(-1289/120,-261/80,153/8,339/20)^T,\notag\\
 Q^{\rm native}&=3161/120,\quad Q^{\rm feedback}=-203/96,\quad
 Q^{\rm red}=11629/480.
 \label{eq:example-result}
\end{align}
The external-wave term of Eq.~\eqref{eq:density-source} is nonzero and is saved in full in the calculation output.

The mixed-quartic script passes 42 exact checks.
These include native density algebra, the 16 independent source components, the generic longitudinal-frequency source \eqref{eq:JAclosed} with arbitrary two transverse polarizations, the actual Hessian inverse, physical--auxiliary orthogonality and the mixed completion-of-square normalization.
Its arbitrary-transverse extension follows from the exact identity \eqref{eq:Rvv}, rather than extrapolation from the longitudinal test.
The finite-solution script passes 14 symbolic checks with \(\alpha,\beta,\kappa,\epsilon,v,w\) kept symbolic.
It reconstructs all Christoffels and all 16 Ricci components directly, then checks Einstein and scalar equations, null gradient, determinant, Ricci scalar, exact amplitude dependence and the rank witness.

The generic mixed formulas have an analytic derivation from the native action.
The four exact realizations check implementation and normalization; they are not advertised as a generic momentum-polynomial proof by sampling.
No new full Compton amplitude, scattering distribution, quartic chart identity, nonlinear Real lift, ghost completion or quantum master equation was verified.

\section{Consequences and the next construction}

The mixed interaction contains two separately necessary pieces: the native minimal scalar contact and the actual quadratic Einstein source propagated through the auxiliary Hessian.
The external-wave correction identifies precisely what is lost by using the on-shell conserved source in an off-shell reduction.
The result is a concrete fourth-order operation on the literal curvature-normalized carrier, retaining its original chart pivots.

The finite construction is a complementary geometric consequence of the nonlinear equations.
A real cubic potential fails its source sign, while a reciprocal transverse coefficient and quartic/quintic potential solve it exactly.
The rank-two scalar Hessian establishes a non-pp coupled sector without asserting an unconstructed Lorentzian interpretation.

The next chart-level calculation is Eq.~\eqref{eq:quartic-defect} with the full quadratic map and native quartics.
Off-shell scalar--graviton mixing also requires pure-graviton quartics for a complete comparison.
Separately, a Lorentzian extension of the finite geometry must add its conjugate-helicity fields while preserving the full scalar stress.
Neither task is completed by cubic cyclic or Real compatibility alone.

\section*{Provenance, chronology and attribution}

The local calculations reported here were completed on 29 October 2026, following a retained linear and cubic Einstein--scalar carrier construction and an on-shell scalar-quartic comparison.
The earlier results are background inputs; the native off-shell mixed contact and exact finite matter-supported sector are the present continuation.
Historical research dates, this local continuation and any future public disclosure are distinct.
No exhaustive originality or priority claim is made.

The operator reports that earlier work was Fable~5-led and that results from July 2026 onward were predominantly produced by GPT models.
DeepSeek and multiple GPT models also participated in cross-review and manuscript preparation.
The present calculations and manuscript were produced by GPT models, including an Astra proposal of the finite real ansatz repair followed by an independent direct tensor verification.
Human authors and affiliations are left blank, and human expert review has not occurred.
The attribution uncertainty is recorded verbatim:
\begin{quote}
I genuinely do not know how authorship, intellectual credit, or ownership should be attributed for these results.
\end{quote}

\bibliographystyle{unsrt}
\bibliography{references}
\end{document}
